% ============================================================
%  Capítulo 4 — Análisis
%  TT 2026-B066 | ESCOM-IPN
%  Frausto Robles Ángel Ali · Rodríguez Verdín Sandoval Vanesa
% ============================================================

\chapter{Análisis}

En este capítulo se definen los requerimientos del sistema a partir de los prototipos
de interfaz desarrollados para el flujo de trabajo del laboratorio colaborador y de
reuniones con el Dr.\ César Augusto Sandino Reyes López, investigador del Laboratorio
de Bioquímica Estructural, Sección de Posgrado, de la ENMyH-IPN. Se documentan
también las reglas de negocio, los riesgos identificados y la factibilidad del proyecto
desde cuatro ángulos: tecnológico, operativo, económico y legal.

Los requerimientos funcionales se organizan por módulo de la interfaz. Los tiempos de
respuesta y criterios de rendimiento se derivan de los atributos de calidad definidos
en la norma ISO/IEC~25010:2023 \cite{iso25010}: \textit{time behaviour},
\textit{functional correctness} y \textit{usability}. El ciclo de levantamiento sigue
los procesos de definición de requerimientos de la norma ISO/IEC/IEEE~12207:2017
\cite{iso12207}.

% -----------------------------------------------------------
\section{Requerimientos}
% -----------------------------------------------------------

\subsection{Requerimientos Funcionales (RF)}

Los requerimientos funcionales están organizados por módulo de la interfaz con tres
niveles de prioridad: \textbf{Alta} (sin esto el prototipo no funciona),
\textbf{Media} (importante pero puede quedar para una iteración posterior) y
\textbf{Baja} (deseable pero no urgente). Cada requerimiento se derivó directamente
de los prototipos de interfaz diseñados \cite{pressman2014,sommerville2011}.

\begin{longtable}{|>{\bfseries}p{1.1cm}|>{\raggedright}p{2.8cm}|>{\raggedright\arraybackslash}p{5.7cm}|p{1.3cm}|p{2.2cm}|}
\caption{Requerimientos Funcionales (RF).
\label{tab:rf}}\\
\hline
\textbf{ID} & \textbf{Módulo} & \textbf{Descripción} & \textbf{Prior.} & \textbf{Actor} \\
\hline
\endfirsthead
\hline
\textbf{ID} & \textbf{Módulo} & \textbf{Descripción} & \textbf{Prior.} & \textbf{Actor} \\
\hline
\endhead
\hline
\endfoot

% ── Autenticación ──────────────────────────────────────────
RF-01 & Autenticación
      & El sistema debe permitir iniciar sesión con correo electrónico y
        contraseña. Las credenciales incorrectas muestran
        un mensaje de error sin revelar cuál campo falló.
      & Alta & Investigador / Admin \\
\hline
RF-02 & Autenticación
      & El sistema debe cerrar la sesión del usuario descartando el token JWT activo
        en el navegador; el servidor no mantiene una lista de tokens cancelados y el
        token tiene una vigencia larga y configurable. La opción de cierre de sesión
        debe ser accesible desde el menú de usuario en todas las pantallas.
      & Alta & Investigador / Admin \\
\hline
RF-03 & Autenticación
      & El sistema debe ofrecer recuperación de contraseña enviando un enlace de
        restablecimiento al correo registrado.
      & Media & Investigador \\
\hline
RF-04 & Autenticación
      & El sistema debe gestionar dos roles con permisos distintos: Investigador
        (puede consultar todos los experimentos del laboratorio) y Administrador
        (mismos permisos que el Investigador, más la capacidad de crear, modificar
        y desactivar cuentas de usuario, y monitorear el estado del sistema).
      & Alta & Sistema \\
\hline

% ── Gestión de usuarios ────────────────────────────────────
RF-05 & Gestión de usuarios
      & El Administrador debe poder crear, modificar y desactivar cuentas de usuario
        desde el panel de administración. La creación requiere identificador institucional,
        nombre, apellidos y correo; el sistema genera una contraseña temporal que el
        usuario cambia al primer ingreso.
      & Alta & Admin \\
\hline
RF-06 & Gestión de usuarios
      & El Investigador debe poder actualizar su nombre, correo y contraseña desde
        su perfil de usuario.
      & Media & Investigador \\
\hline
RF-07 & Gestión de usuarios
      & Solo el Administrador puede crear cuentas de usuario. No existe formulario
        público de autoregistro. El Administrador asigna identificador institucional,
        nombre, apellidos y correo; el sistema genera una contraseña temporal y obliga
        al usuario a cambiarla en el primer ingreso.
      & Alta & Admin \\
\hline

% ── Carga de video ─────────────────────────────────────────
RF-08 & Carga de video
      & El sistema debe aceptar únicamente archivos en formato \texttt{.mp4} o
        \texttt{.mov}. Cualquier otro formato debe ser rechazado con un mensaje de
        error específico antes de intentar subirlo al servidor.
      & Alta & Sistema \\
\hline
RF-09 & Carga de video
      & El sistema debe permitir subir uno o dos videos por tanda: Día~1
        (sesión de estrés, 20~min) y Día~2 (sesión post-tratamiento, 5~min). Ambos
        son opcionales individualmente. El pipeline analiza los últimos 300~segundos
        del video de Día~2 (el resto se descarta); del video de Día~1 analiza
        únicamente los primeros 5~minutos y solo cuando la tanda pertenece al grupo
        control (ver RN-03). Al subir el
        primer video de una tanda nueva, el investigador indica el grupo y el
        tratamiento; si el nombre coincide con un grupo ya registrado en el mismo
        experimento, el sistema sugiere la siguiente tanda y pide confirmación
        antes de guardar.
      & Alta & Investigador \\
\hline
RF-10 & Carga de video
      & El sistema debe validar que el archivo subido sea un video reproducible y que
        esté en horizontal (ancho mayor que alto) antes de encolarlo. Un video grabado
        con la cámara girada (en vertical) no se puede procesar. Si la validación
        falla, lo rechaza con código de error y mensaje claro.
      & Alta & Sistema \\
\hline
RF-11 & Carga de video
      & Al crear un experimento el investigador debe registrar: nombre único del
        experimento, fecha y notas adicionales opcionales. El grupo (control,
        referencia o tratamiento), el tratamiento o condición experimental y el
        número de especímenes de cada tanda (entre~2 y~4) se capturan después,
        por grupo, al subir el primer video de cada uno (ver RF-09).
      & Alta & Investigador \\
\hline
RF-12 & Carga de video
      & El análisis debe iniciarse automáticamente al completarse la carga y
        validación del video, sin intervención del investigador.
      & Alta & Sistema \\
\hline

% ── Análisis conductual ────────────────────────────────────
RF-13 & Análisis conductual
      & El pipeline de análisis debe ejecutar tres etapas secuenciales:
        preprocesamiento del video (decodificación, alineación de la cámara y
        modelo de fondo),
        localización de los cilindros y clasificación de conductas. Cada etapa
        debe tener un indicador de estado visible en la pantalla de progreso.
      & Alta & Sistema \\
\hline
RF-14 & Análisis conductual
      & El sistema debe detectar automáticamente la posición de los cilindros (de dos
        a cuatro) en el encuadre. Si no los encuentra, el pipeline se detiene y
        genera un reporte de diagnóstico.
      & Alta & Sistema \\
\hline
RF-15 & Análisis conductual
      & Mientras el análisis corre, el investigador debe ver el progreso en tiempo
        real con una barra de porcentaje y la etapa activa. El análisis continúa
        en segundo plano aunque el investigador cierre la pestaña.
      & Alta & Investigador \\
\hline
RF-16 & Análisis conductual
      & Si el análisis falla, el sistema debe mostrar la descripción del problema,
        según la etapa en que se detuvo, y un botón para descargar el reporte de
        diagnóstico en PDF. El investigador no puede reiniciar el análisis desde esta
        pantalla (ver RN-04).
      & Alta & Investigador \\
\hline
RF-17 & Análisis conductual
      & El sistema debe clasificar cuadro a cuadro la conducta de cada espécimen: nado
        activo, inmovilidad o escalamiento \cite{armario2021}. Cuando sabe que el
        espécimen no está inmóvil pero no decide entre nado y escalamiento, lo registra
        como \textit{conducta activa}, una cuarta conducta (ver RN-13). El buceo se
        clasifica como nado activo. Las conductas son mutuamente excluyentes (ver
        RN-07).
      & Alta & Sistema \\
      \hline
RF-18 & Análisis conductual
      & El sistema debe asignar una conducta a cada segundo del análisis, sin
        descartar ninguna por su duración, de modo que registre conductas breves (de
        unos 2~segundos) que el análisis manual por bloques de 5~segundos no
        distingue. Para compararse con ese método, los resultados también se
        presentan en bloques de 5~segundos: cada bloque toma la conducta con más
        segundos (si nado activo y escalamiento empatan, es conducta activa) y muestra
        dentro los segundos que lo componen. Los fragmentos en que el espécimen no se
        ve se marcan en el PDF final.
      & Alta & Sistema \\
\hline

% ── Resultados ─────────────────────────────────────────────
RF-19 & Resultados
      & La pantalla de resultados debe mostrar el tiempo total en segundos y el
        porcentaje de cada conducta por espécimen y por sesión, en una tabla
        descargable.
      & Alta & Investigador \\
\hline
RF-20 & Resultados
      & El sistema debe ofrecer un desglose de conductas por minuto para cada espécimen.
      & Media & Investigador \\
\hline
RF-21 & Resultados
      & El sistema debe mostrar una comparación entre los grupos de un experimento
        con el tiempo promedio por conducta de los especímenes de cada grupo,
        calculado con los videos del Día~2.
      & Alta & Investigador \\
\hline
RF-22 & Resultados
      & El investigador debe poder descargar los datos de resultados en formato CSV
        y XLSX desde la pantalla de resultados y desde el dashboard. El archivo debe
        contener las columnas: \texttt{Tiempo (min:seg)}, \texttt{Nado},
        \texttt{Escalamiento}, \texttt{Inmovilidad} y \texttt{Conducta activa}, con
        una fila por minuto
        analizado. El reporte de diagnóstico y el resumen ejecutivo se descargan
        por separado en PDF.
      & Alta & Investigador \\
\hline

% ── Dashboard / gestión de experimentos ──────────────────
RF-23 & Dashboard
      & El dashboard debe mostrar la lista de todos los experimentos del laboratorio
        con nombre, fecha, número de grupos, sesión, estado (en cola, procesando,
        completado, error) y accesos directos a resultados y descarga (ver RN-08).
        La lista es filtrable por estado, fecha o tratamiento; el filtro por
        tratamiento muestra los experimentos que tengan al menos un grupo con
        ese tratamiento.
      & Alta & Investigador \\
\hline
RF-24 & Dashboard
      & El sistema debe mostrar un aviso visible cuando un video está a menos de
        7~días de ser eliminado automáticamente, indicando el nombre del experimento
        y la fecha de vencimiento exacta. El período de 7~días da tiempo suficiente
        para descargar reportes antes del borrado, conforme al atributo
        \textit{user error protection} (protección contra errores del usuario) de
        la ISO/IEC~25010:2023 \cite{iso25010}, que establece que el sistema debe
        proteger al usuario de consecuencias no intencionadas de sus acciones u
        omisiones.
      & Media & Sistema \\
\hline
RF-25 & Dashboard
      & El sistema debe borrar automáticamente los archivos de video 30~días después
        de que el análisis termine exitosamente, notificando al investigador con un
        aviso en la interfaz antes de que ocurra (ver RN-05). Si el disco supera el
        90\,\%, el borrado se adelanta a los videos más antiguos (ver RNF-09).
      & Media & Sistema \\
\hline
RF-26 & Dashboard
      & Los resultados y reportes deben conservarse indefinidamente, aunque el video
        original ya haya sido borrado.
      & Alta & Sistema \\
\hline
RF-27 & Dashboard
      & El investigador debe poder eliminar un experimento propio desde el
        dashboard; el Administrador puede eliminar cualquiera. La acción es
        permanente e irreversible y requiere confirmación explícita (ver RN-09).
      & Media & Investigador / Admin \\
\hline

% ── Administración ─────────────────────────────────────────
RF-28 & Administración
      & El Administrador debe poder ver la lista de todos los usuarios con correo,
        si es administrador o investigador, estado (activo/inactivo)
        y número de experimentos. La lista es filtrable por
        administrador/investigador y estado.
      & Media & Admin \\
\hline
RF-29 & Administración
      & El Administrador debe poder monitorear el estado del sistema: porcentaje de
        uso del disco, espacio disponible en GB, tamaño de la cola de procesamiento,
        número de análisis activos y ETA estimado. El sistema alerta cuando el disco
        supera el 80\,\% de uso.
      & Baja & Admin \\
\hline
RF-30 & Administración
      & El Administrador debe poder ver y gestionar los experimentos de cualquier
        investigador (ver RN-08).
      & Media & Admin \\
\hline
RF-31 & Resultados
      & El sistema debe calcular estadísticas por grupo experimental (control,
        referencia, tratamiento) dentro de cada experimento: media, desviación
        estándar y varianza del tiempo total de cada conducta entre los
        especímenes de ese grupo, agregando todas sus tandas. Estas estadísticas
        deben incluirse en el reporte descargable (CSV/XLSX) junto con los datos
        individuales por espécimen.
      & Media & Investigador \\
\hline
RF-32 & Resultados
      & Cuando un análisis está completado, el sistema debe ofrecer una pantalla de
        revisión segundo a segundo: el video con una línea de tiempo por espécimen,
        coloreada por conducta, donde cada segundo muestra la conducta que puso el
        sistema y el usuario puede corregirla. Cada conducta guarda su origen
        (del sistema, confirmada, corregida o dudosa). Al guardar las correcciones, el sistema recalcula los
        totales por minuto (ver RN-14).
      & Media & Investigador \\
\hline
RF-33 & Notificaciones
      & El sistema debe avisar al usuario correspondiente, en el panel de
        notificaciones de la barra superior y también por correo electrónico, cuando
        un análisis termina, cuando un análisis se detiene por error y, al
        Administrador, cuando el almacenamiento rebasa el 80\,\% (ver RF-29). Los
        correos se envían desde una cuenta de correo configurable en el servidor; si
        el envío falla, el aviso permanece en el panel de notificaciones.
      & Media & Investigador / Admin \\
\hline

\end{longtable}

% -----------------------------------------------------------
\subsection{Requerimientos No Funcionales (RNF)}
% -----------------------------------------------------------

Los tiempos de respuesta se derivan del atributo \textit{time behaviour} de la norma
ISO/IEC~25010:2023 \cite{iso25010}, que establece que un sistema debe responder dentro
de los tiempos que el usuario percibe como aceptables para cada tipo de operación. Para
operaciones interactivas cortas (navegación, filtros, consultas) el umbral estándar es
de 1 a 3 segundos; para operaciones de procesamiento intensivo (análisis de video) el
tiempo aceptable se determina con el usuario durante las pruebas \cite{iso25010}.

\begin{longtable}{|>{\bfseries}p{1.2cm}|>{\raggedright}p{2.9cm}|>{\raggedright}p{5.2cm}|>{\raggedright\arraybackslash}p{3.8cm}|}
\caption{Requerimientos No Funcionales (RNF).
\label{tab:rnf}}\\
\hline
\textbf{ID} & \textbf{Categoría} & \textbf{Descripción} & \textbf{Métrica / Criterio} \\
\hline
\endfirsthead
\hline
\textbf{ID} & \textbf{Categoría} & \textbf{Descripción} & \textbf{Métrica / Criterio} \\
\hline
\endhead
\hline
\endfoot

RNF-01 & Rendimiento
       & Las operaciones de navegación e interacción con la interfaz (cargar el
         dashboard, filtrar experimentos, abrir resultados) deben responder en menos
         de 3~s. Derivado del atributo \textit{time behaviour} de la ISO/IEC~25010:2023
         \cite{iso25010}.
       & Tiempo de respuesta $\leq 3$~s en el 95\,\% de las solicitudes en condiciones
         normales de red. \\
\hline
RNF-02 & Calidad del análisis
       & El tiempo de procesamiento no es el criterio prioritario para el laboratorio.
         El criterio de calidad es el desempeño del clasificador por conducta: se considera
         aceptable un F1 $\geq 85$\,\% por conducta. Este umbral corresponde a la
         variabilidad aceptada entre analistas humanos, declarada por el investigador
         colaborador en la entrevista de levantamiento (véase la
         justificación de umbrales en la sección~\ref{sec:umbrales}).
       & F1 $\geq 85$\,\% por conducta (nado, inmovilidad, escalamiento) sobre el
         conjunto de prueba de la fase de validación (TT2). \\
\hline
RNF-03 & Disponibilidad
       & El sistema debe estar disponible las 24 horas del día, los 7 días de la
         semana, para permitir que los investigadores dejen análisis corriendo fuera
         del horario laboral.
       & $\geq 95$\,\% de disponibilidad mensual medida en el entorno de despliegue. \\
\hline
RNF-04 & Seguridad
       & Solo los usuarios autenticados pueden acceder al sistema. La autenticación
         usa tokens JWT con tiempo de expiración configurable. Todos los
         investigadores con cuenta activa pueden consultar los experimentos y
         resultados de cualquier otro investigador del laboratorio (ver RN-08).
       & Autenticación con JWT verificada en pruebas de integración. Acceso a datos
         de otros investigadores restringido a usuarios con sesión activa. \\
\hline
RNF-05 & Seguridad
       & Las contraseñas deben almacenarse cifradas con sal, nunca en texto plano.
         Las contraseñas temporales asignadas por el Admin se marcan para cambio
         obligatorio al primer ingreso.
       & Hash con sal usando bcrypt. Política de contraseña temporal verificada
         en prueba funcional. \\
\hline
RNF-06 & Usabilidad
       & Un investigador sin experiencia en programación debe poder crear un
         experimento, subir un video, esperar el análisis y descargar el reporte
         sin necesitar ayuda técnica. Derivado del atributo \textit{usability} de
         la ISO/IEC~25010:2023 \cite{iso25010}.
       & Prueba con al menos un usuario real del laboratorio durante la fase de
         validación y robustez. Tasa de completación de tarea $\geq$ 90\% sin
         asistencia. \\
\hline
RNF-07 & Compatibilidad
       & El sistema debe funcionar desde Chrome, Firefox y Edge en sus versiones
         actuales, sin instalar nada en el equipo del usuario.
       & Pruebas en los tres navegadores en sus versiones actuales al cierre de
         la fase de implementación de la interfaz web. \\
\hline
RNF-08 & Almacenamiento
       & Los videos se borran automáticamente 30~días después de que el análisis
         termina exitosamente. Los resultados y reportes se conservan
         indefinidamente. El sistema alerta al usuario cuando un video está a menos
         de 7~días de ser borrado.
       & Tarea programada (\textit{cron job}) diaria. Aviso visible en dashboard
         con indicador de días restantes. \\
\hline
RNF-09 & Almacenamiento
       & El Administrador recibe una alerta en el panel cuando el uso del disco
         supera el 80\,\%. A los 90\,\%, el sistema elimina automáticamente los
         videos más antiguos, notificando al investigador correspondiente.
       & Umbral de alerta: 80\,\%. Umbral de limpieza automática: 90\,\%.
         Los valores 80/90\,\% son umbrales estándar de gestión de capacidad
         en sistemas de almacenamiento \cite{pressman2014}; dejan margen
         operativo antes de que el disco se llene y evitan pérdida de datos
         por falta de espacio. \\
\hline
RNF-10 & Portabilidad
       & El sistema debe poder levantarse con Docker Compose en el servidor
         institucional provisto por la ESCOM, sin dependencias adicionales en el
         sistema anfitrión. Derivado del proceso de despliegue de la
         ISO/IEC/IEEE~12207:2017 \cite{iso12207}.
       & Despliegue documentado y verificado en al menos dos entornos distintos. \\
\hline
RNF-11 & Escalabilidad
       & El módulo de clasificación debe poder reconfigurarse para soportar
         paradigmas conductuales distintos al FST sin modificar la plataforma web.
       & Diseño modular verificado en la revisión de arquitectura al cierre
         del bloque de diseño. \\
\hline
RNF-12 & Mantenibilidad
       & El código debe estar documentado y separado en módulos independientes
         (backend, pipeline de análisis de video, frontend), siguiendo PEP~8 en Python y ESLint
         en JavaScript.
       & Revisión de calidad de código en cada cierre de bloque de
         implementación. \\
\hline
RNF-13 & Concurrencia
       & El sistema debe soportar hasta 4 usuarios registrados trabajando en el
         mismo periodo de tiempo. El uso simultáneo estricto no se requiere: los
         análisis se encolan y procesan de forma secuencial (ver RN-04). Este valor
         fue declarado por el investigador colaborador en la entrevista de
         levantamiento (véase la justificación de umbrales en la
         sección~\ref{sec:umbrales}).
       & Pruebas con 4 sesiones de usuario activas simultáneamente y al menos 1
         análisis en cola, sin degradación de respuesta $> 3$~s en operaciones
         de interfaz. \\
\hline

\end{longtable}

% -----------------------------------------------------------
\subsection{Justificación de umbrales y valores numéricos}
\label{sec:umbrales}
% -----------------------------------------------------------

Los requerimientos y reglas de negocio contienen valores numéricos concretos.
A continuación se documenta el origen de cada uno para que puedan auditarse y
ajustarse en la fase de validación (TT2).

\begin{itemize}
  \item \textbf{F1 mínimo por conducta: $\geq 85$\,\%.}
        En la investigación conductual con roedores, la práctica estándar para
        validar un sistema automatizado es comparar su clasificación contra la
        anotación manual de analistas entrenados y exigir que la concordancia
        alcance o supere el nivel de acuerdo típico entre dos evaluadores humanos
        \cite{slattery2012,della_valle2025,nandi2021_dbscorer}. Ese nivel de
        acuerdo entre humanos oscila entre el 80 y el 90\,\% en la codificación
        de conductas del FST, dependiendo de la experiencia del evaluador y de
        la nitidez del video \cite{slattery2012,della_valle2025}. El Dr.\ César
        Augusto Sandino Reyes López declaró en la entrevista de levantamiento que
        la variabilidad aceptada en su laboratorio es del 85\,\%, valor que se
        ubica en ese rango y que es consistente con los umbrales reportados en
        sistemas automatizados comparables \cite{nandi2021_dbscorer,meng2024_dsaan}.
        El F1 por conducta se evaluará sobre el conjunto de prueba de la fase de
        validación en TT2.

  \item \textbf{Sin duración mínima de episodio; bloques de 5~segundos como referencia.}
        En la entrevista de levantamiento, el investigador colaborador indicó que en
        el análisis manual una conducta se contabiliza si se mantiene alrededor de 3
        a 5~segundos. El sistema no adopta ese mínimo: asigna una conducta a cada
        segundo y detecta conductas breves de unos 2~segundos, que es justo lo que el
        análisis a ojo no alcanza a registrar. El criterio manual se conserva como
        referencia: los resultados también se agrupan en bloques de 5~segundos, con la
        conducta que más segundos ocupó en cada bloque, para compararlos con el
        análisis manual (ver RF-18).

  \item \textbf{Período de retención de videos: 30~días.}
        Confirmado por el investigador colaborador en la entrevista de levantamiento
        como plazo suficiente para descargar reportes antes del borrado automático
        (ver RN-05, RF-24, RNF-08).

  \item \textbf{Capacidad de usuarios: 4 personas.}
        Declarado por el investigador colaborador: el laboratorio trabaja con un
        máximo de tres tesistas más el responsable del laboratorio, todos usando
        el sistema en el mismo periodo académico aunque no de forma simultánea
        (ver RNF-13).

  \item \textbf{Umbral de alerta de disco: 80\,\% / limpieza: 90\,\%.}
        Umbrales estándar de gestión de capacidad en sistemas de almacenamiento
        \cite{pressman2014}. Dejan margen operativo suficiente antes de que el
        disco se llene y evitan pérdida de datos por falta de espacio. Se
        revisarán según el volumen real de videos generados por semestre (ver RNF-09).
\end{itemize}

% -----------------------------------------------------------
\subsection{Reglas de negocio}
% -----------------------------------------------------------

Las reglas de negocio son los límites que el sistema siempre debe respetar,
independientemente de cómo esté implementado. Algunas vienen del protocolo
experimental del FST, otras del diseño de la interfaz y otras del contexto
institucional del laboratorio \cite{pressman2014}. Los requerimientos que dependen
de estas reglas las referencian explícitamente.

\begin{longtable}{|>{\bfseries}p{1.1cm}|>{\raggedright}p{8.1cm}|>{\raggedright\arraybackslash}p{3.9cm}|}
\caption{Reglas de negocio.
\label{tab:rn}}\\
\hline
\textbf{ID} & \textbf{Regla} & \textbf{Justificación / RF relacionados} \\
\hline
\endfirsthead
\hline
\textbf{ID} & \textbf{Regla} & \textbf{Justificación / RF relacionados} \\
\hline
\endhead
\hline
\endfoot

RN-01 & Las credenciales incorrectas muestran un mensaje de error genérico que no
        especifica si el correo o la contraseña es lo que falló.
      & Seguridad básica contra enumeración de usuarios. Ver RF-01. \\
\hline
RN-02 & Un video pertenece a exactamente una tanda (y, a través de ella, a un
        grupo y un experimento). No se puede reutilizar un video ya subido en
        una tanda distinta.
      & Trazabilidad de resultados. Ver RF-09. \\
\hline
RN-03 & Una tanda tiene como máximo dos videos: uno del Día~1 (20~min) y uno del
        Día~2 (5~min). Un experimento, al tener varios grupos y tandas, puede
        acumular muchos más de dos videos en total.
      & Así funciona el protocolo FST del laboratorio \cite{slattery2012}. Ver RF-09. \\
\hline
RN-04 & El análisis lo inicia el sistema automáticamente al terminar la carga. El
        investigador no puede iniciarlo, pausarlo, cancelarlo ni reiniciarlo desde
        la interfaz.
      & Evita estados inconsistentes en la cola. Ver RF-12, RF-16. \\
\hline
RN-05 & Los videos se borran automáticamente 30~días después de que el análisis
        termina exitosamente. La acción es irreversible. El sistema avisa con al
        menos 7~días de anticipación. Excepción: si el disco supera el 90\,\%, se
        adelanta el borrado de los videos más antiguos aunque no hayan cumplido
        los 30~días y se notifica al investigador de cada experimento afectado
        (ver RNF-09).
      & Control de almacenamiento. Ver RF-24, RF-23, RNF-08. \\
\hline
RN-06 & Los resultados y reportes (PDF y CSV) se conservan indefinidamente aunque
        el video original ya haya sido borrado.
      & Los datos de investigación tienen valor a largo plazo. Ver RF-26. \\
\hline
RN-07 & En cada cuadro del video, a cada espécimen se le asigna exactamente una
        conducta: nado activo, inmovilidad, escalamiento o conducta activa. No pueden
        coexistir dos conductas al mismo tiempo para el mismo espécimen.
      & El FST define estas conductas como mutuamente excluyentes
        \cite{armario2021,cryan2005}. Ver RF-17, RF-18. \\
\hline
RN-08 & Todos los investigadores con cuenta activa pueden consultar los experimentos
        y resultados de cualquier otro investigador del laboratorio, pero solo pueden
        eliminar los propios. El Administrador puede además eliminar cualquier
        experimento, y modificar y desactivar cuentas.
      & Definido en entrevista de levantamiento. Ver RF-04, RNF-04, RF-30. \\
\hline
RN-09 & Borrar un experimento es permanente e irreversible: se eliminan el video,
        los resultados y los reportes. El sistema requiere confirmación explícita
        antes de ejecutarlo.
      & Previene borrados accidentales. Ver RF-27. \\
\hline
RN-10 & El sistema rechaza cualquier archivo que no sea \texttt{.mp4} o
        \texttt{.mov}, y cualquier video en vertical (grabado con la cámara girada),
        antes de intentar subirlo al servidor. La validación de formato ocurre en el
        cliente. Un video girado no se puede procesar.
      & Evita consumir recursos con archivos incompatibles. Los videos reales del
        laboratorio son \texttt{.mov}. Ver RF-08, RF-10. \\
\hline
RN-11 & Si el video no se puede abrir o el pipeline no encuentra los cilindros, se
        detiene y el sistema genera automáticamente un reporte de diagnóstico
        descargable que indica el motivo del fallo. Los videos en que sí se
        encuentran continúan el análisis de forma normal.
      & Evita que el sistema produzca resultados incorrectos sin advertir al
        investigador. Véase la sección~\ref{sec:umbrales}. Ver RF-14, RF-16. \\
\hline
RN-12 & La comparación entre grupos usa solo los videos del Día~2. Los
        especímenes cuyo video del Día~2 aún no está analizado se señalan como
        faltantes.
      & El Día~1 solo verifica que los especímenes lleguen con el mismo nivel de
        estrés; el efecto del tratamiento se mide en el Día~2. Ver RF-21. \\
\hline
RN-13 & Cuando el clasificador sabe que el espécimen no está inmóvil pero no
        decide entre nado activo y escalamiento, esos segundos se registran como
        \textit{conducta activa}, una cuarta conducta. Un mismo análisis puede
        mezclar las cuatro conductas, y los reportes muestran la conducta activa
        como una columna más, de modo que el investigador pueda interpretar los
        resultados en consecuencia.
      & Contemplado en el alcance (cap.~1) y en el riesgo R-08. Corresponde al
        nivel de Porsolt (inmóvil o no inmóvil) cuando no se separa lo activo en
        nado y escalamiento. Ver RF-17, RF-18. \\
\hline
RN-14 & Solo se puede revisar un análisis completado. Al revisar, el usuario corrige
        las etiquetas segundo a segundo; cada etiqueta guarda su origen. Al guardar, el sistema
        recalcula los totales por minuto de los especímenes corregidos y los
        reportes guardados del experimento se regeneran con los datos nuevos.
      & La revisión humana corrige al clasificador sin sustituirlo. Ver RF-32. \\
\hline

\end{longtable}

% -----------------------------------------------------------
\section{Análisis de las herramientas a utilizar}
% -----------------------------------------------------------

La elección del \textit{stack} tecnológico se hizo evaluando madurez, compatibilidad,
disponibilidad de documentación y experiencia previa del equipo. Las herramientas se
agrupan en cuatro áreas.

\subsection{Visión por computadora}

\begin{longtable}{|p{2.7cm}|p{4.4cm}|p{4.8cm}|p{1.3cm}|}
\caption{Técnicas de visión por computadora.
\label{tab:herr_cv}}\\
\hline
\textbf{Herramienta} & \textbf{Características principales} & \textbf{Por qué se eligió} & \textbf{GPU req.} \\
\hline
\endfirsthead
\hline
\textbf{Herramienta} & \textbf{Características principales} & \textbf{Por qué se eligió} & \textbf{GPU req.} \\
\hline
\endhead
\hline
\endfoot

OpenCV 4.x \cite{bradski2000}
    & Biblioteca estándar para procesamiento de imágenes y video en Python. Incluye
      detección de puntos característicos, alineación de imágenes, sustracción de
      fondo, morfología y regiones conexas.
    & Es la herramienta de referencia para este tipo de trabajo y es compatible con
      PyTorch sin configuración extra.
    & No \\
\hline
Sustracción de fondo por percentil \cite{piccardi2004}
    & El fondo se estima como el percentil 10 de cada píxel a lo largo del video,
      con los cuadros ya alineados. Como el espécimen es claro sobre agua oscura,
      un percentil bajo lo borra del fondo.
    & Más simple y predecible que el GMM \cite{zivkovic2004} para el caso de cámara
      fija. Se prefirió a la mediana porque esta deja incrustado en el fondo a un
      espécimen que pasa mucho tiempo inmóvil, y luego no se le puede separar.
    & No \\
\hline

\end{longtable}

\subsection{Aprendizaje profundo para clasificación de conductas}

\begin{longtable}{|p{2.7cm}|p{5.2cm}|p{5.2cm}|}
\caption{Arquitecturas de aprendizaje profundo para clasificación de conductas.
\label{tab:herr_dl}}\\
\hline
\textbf{Herramienta / Arquitectura} & \textbf{Características principales} & \textbf{Por qué se eligió} \\
\hline
\endfirsthead
\hline
\textbf{Herramienta / Arquitectura} & \textbf{Características principales} & \textbf{Por qué se eligió} \\
\hline
\endhead
\hline
\endfoot

R(2+1)D-18 \cite{tran2018}
    & Red convolucional espacio-temporal de 18~capas con conexiones residuales
      \cite{he2016}: una ResNet-18 para video. Recibe clips de 16 fotogramas y se
      preentrena con Kinetics-400 (videos de acciones humanas) para aplicar
      transfer learning \cite{pan2010}.
    & Candidata para el clasificador de conducta. Se compara contra el
      clasificador sobre rasgos; la elección se hace en TT-II con F1 por
      conducta. \\
\hline
PyTorch \cite{paszke2019}
    & Framework de aprendizaje profundo para entrenamiento e inferencia de redes
      neuronales. Su biblioteca torchvision incluye la R(2+1)D-18 preentrenada.
    & Es el ecosistema de referencia para visión por computadora en Python.
      Compatible con Docker para el despliegue. \\
\hline
scikit-learn \cite{pedregosa2011}
    & Biblioteca de ML con clasificadores de \textit{gradient boosting}, métricas y
      validación cruzada.
    & Implementa el clasificador sobre rasgos (\texttt{HistGradientBoostingClassifier}),
      la otra candidata para el clasificador de conducta, y calcula las métricas de
      clasificación (F1, precisión y \textit{recall} por conducta) y la validación
      cruzada. \\
\hline
Etiquetador del equipo
    & Herramienta propia que propone la conducta de cada espécimen por bloque y por
      segundo, y permite revisarla y corregirla en el navegador. Guarda las hojas
      corregidas en formato CSV, compatible con Python.
    & Los autores del proyecto revisan y corrigen sus propuestas, sin usar registros
      conductuales del laboratorio como referencia, y generan así el dataset de
      entrenamiento del clasificador. \\
\hline

\end{longtable}

\subsection{Desarrollo web e infraestructura}

\begin{longtable}{|p{2.7cm}|p{5.2cm}|p{5.2cm}|}
\caption{Herramientas de desarrollo web e infraestructura.
\label{tab:herr_web}}\\
\hline
\textbf{Herramienta} & \textbf{Características principales} & \textbf{Por qué se eligió} \\
\hline
\endfirsthead
\hline
\textbf{Herramienta} & \textbf{Características principales} & \textbf{Por qué se eligió} \\
\hline
\endhead
\hline
\endfoot

Flask (Python) \cite{grinberg2018}
    & Microframework ligero para APIs REST. Extensiones para JWT
      (Flask-JWT-Extended), ORM (SQLAlchemy) y tareas asíncronas.
    & Suficiente para el tamaño del prototipo. Integra el pipeline de análisis de video sin
      necesitar un framework más pesado. \\
\hline
React.js + Vite
    & Framework frontend basado en componentes. Vite reduce los tiempos de
      compilación durante el desarrollo.
    & Adecuado para la barra de progreso en tiempo real y la interfaz reactiva
      que requiere el sistema. \\
\hline
PostgreSQL
    & Base de datos relacional con transacciones ACID y código abierto.
    & Cubre los patrones de acceso del sistema y se integra con SQLAlchemy. \\
\hline
Docker + Docker Compose \cite{merkel2014}
    & Empaqueta todos los servicios (backend, frontend, base de datos, worker) en
      contenedores y los levanta con un solo comando.
    & Cumple el atributo de portabilidad de la ISO/IEC/IEEE~12207:2017
      \cite{iso12207}. El laboratorio no instala nada más allá de Docker. \\
\hline

\end{longtable}

% -----------------------------------------------------------
\section{Análisis de riesgos}
% -----------------------------------------------------------

La identificación de riesgos sigue el modelo de la norma ISO~31000 \cite{iso31000}, que
define el riesgo como la combinación entre la probabilidad de un evento y su impacto.
Los riesgos se clasifican con un semáforo de tres niveles y se ordenan de mayor a
menor criticidad: \textcolor{red!70!black}{\textbf{alto}} (intervención inmediata),
\textcolor{orange!80!black}{\textbf{medio}} (monitoreo continuo) y
\textcolor{green!55!black}{\textbf{bajo}} (seguimiento periódico).
Para cada riesgo se documenta una acción concreta de mitigación.

\begin{longtable}{|>{\bfseries}p{0.9cm}|p{2.0cm}|p{4.5cm}|p{5.9cm}|}
\caption{Matriz de riesgos del proyecto.
\label{tab:riesgos}}\\
\hline
\textbf{ID} & \textbf{Nivel} & \textbf{Descripción} & \textbf{Mitigación} \\
\hline
\endfirsthead
\hline
\textbf{ID} & \textbf{Nivel} & \textbf{Descripción} & \textbf{Mitigación} \\
\hline
\endhead
\hline
\endfoot

% ---- NIVEL ALTO (rojo) ----
R-01 & \cellcolor{red!12}\textcolor{red!70!black}{\Large$\bullet$}~\textbf{Alto}
     & Videos con reflejos, iluminación variable o especímenes de pelaje claro que
       hacen fallar el pipeline (no se encuentran los cilindros, ver RN-11).
     & Definir con el laboratorio los requisitos mínimos de calidad de video.
       El sistema deduce el fondo y la posición de los cilindros de cada video, sin
       parámetros fijados a mano para un encuadre; si aun así no se encuentran los
       cilindros, se genera el reporte de diagnóstico (ver RN-11). \\
\hline
R-02 & \cellcolor{red!12}\textcolor{red!70!black}{\Large$\bullet$}~\textbf{Alto}
     & El dataset de cuadros etiquetados no alcanza para entrenar un clasificador
       con el F1 objetivo.
     & Ampliar el dataset anotando más cuadros con el etiquetador del equipo y aplicar
       \textit{data augmentation}. Si el escalamiento no se distingue del nado
       con la precisión objetivo, los segundos en que no se distinguen se registran
       como ``conducta activa'' (ver RN-13). \\
\hline

% ---- NIVEL MEDIO (amarillo) ----
R-03 & \cellcolor{orange!15}\textcolor{orange!80!black}{\Large$\bullet$}~\textbf{Medio}
     & El laboratorio no proporciona los registros conductuales en el tiempo
       necesario para el entrenamiento y validación del clasificador.
     & Acordar un calendario de entrega de registros con el laboratorio al inicio
       de la fase de implementación y preparar los materiales de anotación con
       anticipación para aprovechar cada sesión con el Dr.\ Sandino. \\
\hline
R-04 & \cellcolor{orange!15}\textcolor{orange!80!black}{\Large$\bullet$}~\textbf{Medio}
     & El tiempo de análisis de un video de 5~min supera lo que el laboratorio
       puede asumir en su flujo de trabajo.
     & Medir tiempos por etapa del pipeline. Evaluar reducción de resolución o
       uso de GPU. El umbral aceptable de tiempo se define con el laboratorio
       durante la fase de pruebas (ver RNF-02). \\
\hline
R-05 & \cellcolor{orange!15}\textcolor{orange!80!black}{\Large$\bullet$}~\textbf{Medio}
     & La integración del módulo de visión por computadora con el \textit{backend}
       vía cola asíncrona resulta más compleja de lo esperado.
     & Definir contratos de \textit{API} al inicio de la fase de implementación
       y ejecutar una prueba de integración de extremo a extremo antes de agregar
       más funcionalidad. \\
\hline
R-07 & \cellcolor{orange!15}\textcolor{orange!80!black}{\Large$\bullet$}~\textbf{Medio}
     & El laboratorio solicita cambios en los requerimientos a mitad del desarrollo.
     & Diseño modular del clasificador para incorporar clases adicionales sin
       reescribir la arquitectura. Gestionar los cambios dentro del proceso de
       revisión de \textit{sprint} conforme a la metodología Scrum. \\
\hline
R-08 & \cellcolor{orange!15}\textcolor{orange!80!black}{\Large$\bullet$}~\textbf{Medio}
     & El escalamiento y el nado activo se ven tan parecidos desde la vista lateral
       que el modelo no los distingue con suficiente precisión.
     & Usar en el entrenamiento únicamente clips inequívocos de escalamiento
       (movimiento de patas delanteras hacia las paredes del cilindro). Si no se
       alcanza la precisión objetivo, registrar como ``conducta activa'' los
       segundos en que no se distinguen (ver RN-13). \\
\hline

% ---- NIVEL BAJO (verde) ----
R-06 & \cellcolor{green!12}\textcolor{green!55!black}{\Large$\bullet$}~\textbf{Bajo}
     & El disco del servidor se llena por acumulación de videos durante el período
       de retención de 30~días.
     & El sistema genera alertas al 80\,\% de uso de disco y ejecuta limpieza
       automática al 90\,\% (ver RF-25, RF-29, RNF-09). \\
\hline

\end{longtable}

% -----------------------------------------------------------
\section{Análisis de factibilidad}
% -----------------------------------------------------------

\subsection{Factibilidad tecnológica}
\label{sec:fact_tecnologica}

La factibilidad tecnológica consiste en evaluar si el proyecto puede desarrollarse con
la infraestructura técnica, el conocimiento especializado y los recursos disponibles.
El objetivo es determinar si las tecnologías actuales y el personal son suficientes para
implementar el prototipo de forma eficiente y sostenible \cite{pressman2014}.

\subsubsection{Evaluación de recursos actuales y brechas tecnológicas}

El equipo de desarrollo cuenta con experiencia en las herramientas del stack tecnológico
seleccionado ---Python, OpenCV, PyTorch, Flask y React--- adquirida en cursos de la
carrera en la ESCOM. Todas las herramientas son de código abierto, están bien
documentadas y tienen comunidades activas que reducen el riesgo de incompatibilidades.

La brecha tecnológica más relevante identificada es el dataset de cuadros etiquetados
para entrenar el clasificador de conducta. Con el número de videos disponibles actualmente
en el laboratorio, ese dataset requiere un esfuerzo de anotación manual con el etiquetador del equipo que
es viable pero no trivial. La estrategia de transfer learning con una red preentrenada
reduce ese requerimiento a unos pocos cientos de ejemplos por clase, lo que está dentro
del alcance del proyecto \cite{pan2010,meng2024_dsaan}. El clasificador sobre rasgos
corre en procesador, sin GPU. Para entrenar la red convolucional, si el laboratorio no
dispone de GPU, Google Colab ofrece hasta 12~horas de GPU por sesión de forma gratuita,
lo que es suficiente para su ajuste fino.

En cuanto al hardware para el despliegue, el sistema no requiere infraestructura
especializada: cualquier equipo con Docker instalado puede correr el prototipo, lo que
elimina dependencias de hardware específico en el laboratorio.

\subsubsection{Simulación, modelado y pruebas}

Antes de la implementación completa del sistema, el pipeline de análisis se probó con
videos del laboratorio: localiza de forma automática los cilindros y la línea de agua
en grabaciones de distintos días y encuadres, con vista lateral y de dos a cuatro
cilindros simultáneos. Esa prueba reduce el riesgo técnico de la etapa de localización
de los cilindros antes de comprometer recursos en el desarrollo completo.

Para el clasificador de conducta, estudios previos reportan precisiones superiores al
88\,\% con datasets de menos de 100 imágenes cuando se usa transfer learning desde
ImageNet \cite{meng2024_dsaan}, lo que da respaldo empírico a la viabilidad de esa
etapa con el dataset disponible.

\subsubsection{Evaluación de compatibilidad e integración}

El prototipo tiene una arquitectura modular que separa el pipeline de análisis de video
(Python, OpenCV, scikit-learn y PyTorch) de la plataforma web (Flask, React, PostgreSQL) mediante una
cola de tareas asíncrona. Esa separación permite desarrollar y probar cada módulo de
forma independiente y facilita la integración progresiva conforme avanza el desarrollo.

La contenedorización con Docker garantiza que el sistema se comporta igual en el equipo
de desarrollo, en el laboratorio y en cualquier servidor institucional, lo que elimina
la clase más común de problemas de compatibilidad en el despliegue \cite{merkel2014}.

\subsection{Factibilidad operativa}

La factibilidad operativa analiza si el sistema encajará en la rutina de trabajo real
del laboratorio una vez entregado, es decir, si los usuarios lo adoptarán sin que
represente una carga adicional sobre la que ya tienen. Tres preguntas concretas guían
este análisis: ¿quién lo usará y qué tan familiarizado está con herramientas digitales?,
¿cómo se integra con el flujo de trabajo actual del laboratorio?, y ¿qué pasa si algo
falla durante un experimento?

\subsubsection{Perfil del usuario y contexto de uso}

El usuario principal es el Dr.\ César Augusto Sandino Reyes López, investigador del
Laboratorio de Bioquímica Estructural (Sección de Posgrado, ENMyH-IPN), junto con los
estudiantes de posgrado que participan en los experimentos. Ninguno de ellos es
desarrollador de software, lo que impone una restricción de diseño clara: la interfaz
debe requerir cero configuración técnica del lado del usuario. Cargar un video,
asignar grupos y descargar el reporte deben ser las únicas acciones necesarias.

El laboratorio realiza experimentos de nado forzado (FST) sin un calendario fijo; en
un semestre de actividad regular son del orden de 3 a 5, cada uno con tres grupos (control, referencia y tratamiento experimental) de entre
6 y 8 especímenes cada uno. Cada experimento genera entre 6 y 9 videos de 5~min en
el segundo día del protocolo (uno por tanda, con dos o tres tandas por grupo). El análisis de esos videos es el cuello de botella actual:
anotarlos manualmente ocupa entre 1.5 y 2~h cada uno (con sus dos a cuatro
especímenes), lo que convierte el análisis en
una actividad que compite con el tiempo de investigación del laboratorio.

\subsubsection{Integración con el protocolo experimental}

El sistema no modifica el protocolo del laboratorio; se inserta al final, cuando los
videos ya están grabados. El investigador sube los archivos desde cualquier computadora
con navegador ---sin instalar nada--- y el sistema los procesa de forma asíncrona
mientras él continúa con otras tareas. Los resultados se descargan en los formatos que
el laboratorio ya usa para sus análisis estadísticos: CSV y XLSX para datos tabulares
y PDF para el reporte diagnóstico de calidad del video.

Esta ausencia de cambio en el protocolo de grabación es una ventaja operativa
importante: el laboratorio no necesita adaptar su metodología ni adquirir equipo nuevo.
Los videos grabados con la cámara web actual son suficientes, siempre que el
cilindro sea visible en su totalidad desde la vista lateral.

\subsubsection{Criterios de aceptación operativa}

Con base en las reuniones sostenidas con el Dr.\ Sandino durante el arranque del
proyecto, los criterios que determinan si el sistema es operativamente aceptable son:

\begin{itemize}
  \item El tiempo de procesamiento de un video de 5~min no debe superar el tiempo
        que tomaría anotarlo manualmente (entre 1.5 y 2~h). Cualquier resultado en
        ese rango o por debajo ya representa una ganancia operativa.
  \item El reporte generado debe incluir, al menos, la duración acumulada por conducta
        (nado activo, inmovilidad, escalamiento y conducta activa) desglosada por
        minuto, de modo que sea directamente comparable con la tabla que el
        investigador construye hoy a mano.
  \item Cuando el video no se pueda procesar (no se abre o no se encuentran los
        cilindros, según RN-11), el sistema debe notificarlo explícitamente en el
        reporte de diagnóstico en lugar de entregar un resultado silencioso
        potencialmente erróneo.
\end{itemize}

Estos tres criterios se verificarán durante la etapa de validación en TT-II con los
videos del laboratorio y las métricas de clasificación definidas en RNF-02.

\subsection{Factibilidad económica}

La factibilidad económica determina si el proyecto puede concluirse dentro de los
recursos con que cuenta el equipo de desarrollo. El principal costo del proyecto es
el recurso humano: dos integrantes que destinan aproximadamente 500~horas cada uno
durante TT-I y TT-II combinados. Los demás costos corresponden a servicios de cómputo
en la nube, documentación e insumos de soporte. Dado que las herramientas del
\textit{stack} son de código abierto y el laboratorio ya dispone de equipo de
grabación, no se requiere adquisición de hardware especializado.

\subsubsection{Costos fijos}

Se consideran costos fijos aquellos que no dependen del ritmo de avance ni del número
de iteraciones del desarrollo:

\begin{longtable}{|p{6.5cm}|p{6.8cm}|}
\caption{Costos fijos estimados.
\label{tab:costos_fijos}}\\
\hline
\textbf{Rubro} & \textbf{Monto aproximado (MXN)} \\
\hline
\endfirsthead
\hline
\textbf{Rubro} & \textbf{Monto aproximado (MXN)} \\
\hline
\endhead
\hline
\endfoot

Recurso humano --- 2 desarrolladores $\times$ 500~h~c/u a \$40/h
      (tarifa referencial de pasante de ingeniería)
    & \$40{,}000 \\
\hline
Equipo de cómputo y cámara de grabación
    & \$0 --- el Laboratorio de Bioquímica Estructural (ENMyH-IPN) pone a disposición
      el equipo existente. \\
\hline
Cómputo acelerado por GPU --- Google Colab Pro (plan anual, entrenamiento
      del clasificador)
    & \$3{,}300 \\
\hline
Conectividad e infraestructura (internet y electricidad durante el desarrollo,
      12~meses)
    & \$4{,}500 \\
\hline
Herramientas de desarrollo (Python, Flask, React, PostgreSQL, Docker, OpenCV,
      scikit-learn, PyTorch)
    & \$0 --- distribución bajo licencias de código abierto sin costo para uso
      académico. \\
\hline
Impresión y encuadernación de dos ejemplares del documento recepcional
    & \$2{,}400 \\
\hline
Trámites y gestiones administrativas (titulación, registro)
    & \$800 \\
\hline
\textbf{Subtotal costos fijos} & \textbf{\$51{,}000} \\
\hline

\end{longtable}

\subsubsection{Costos variables}

Los costos variables están condicionados al avance del proyecto y solo se incurren si
determinadas situaciones se materializan durante el desarrollo:

\begin{longtable}{|p{6.5cm}|p{6.8cm}|}
\caption{Costos variables estimados.
\label{tab:costos_variables}}\\
\hline
\textbf{Rubro} & \textbf{Monto aproximado (MXN)} \\
\hline
\endfirsthead
\hline
\textbf{Rubro} & \textbf{Monto aproximado (MXN)} \\
\hline
\endhead
\hline
\endfoot

Cómputo adicional en la nube (almacenamiento ampliado en Google Drive o similar,
      solo si el dataset de video supera la capacidad del nivel gratuito)
    & \$0--\$1{,}500 \\
\hline
Insumos de soporte audiovisual (almacenamiento externo, adaptadores, cables)
    & \$2{,}000 \\
\hline
Formación técnica complementaria (acceso a cursos especializados de visión por
      computadora y despliegue con Docker)
    & \$2{,}000 \\
\hline
\textbf{Subtotal costos variables} & \textbf{\$5{,}500} \\
\hline

\end{longtable}

\subsubsection{Resumen y conclusión económica}

\begin{longtable}{|p{6.5cm}|p{6.8cm}|}
\caption{Estimación total de costos del proyecto.
\label{tab:costos_total}}\\
\hline
\textbf{Concepto} & \textbf{Monto aproximado (MXN)} \\
\hline
\endfirsthead
\hline
\textbf{Concepto} & \textbf{Monto aproximado (MXN)} \\
\hline
\endhead
\hline
\endfoot

Costos fijos & \$51{,}000 \\
\hline
Costos variables & \$5{,}500 \\
\hline
\textbf{Costo total estimado} & \textbf{\$56{,}500} \\
\hline

\end{longtable}

Con un costo total estimado de \$56{,}500~MXN, el rubro principal es el recurso
humano (\$40{,}000), que refleja el esfuerzo real de desarrollo en un contexto
académico. El resto de la inversión corresponde a servicios de cómputo en la nube para
el entrenamiento del clasificador, conectividad, documentación e insumos menores. El
retorno es inmediato para el laboratorio: anotar manualmente un video de 5~min toma
entre 1.5 y 2~h; con alrededor de 32 a 40~videos en un semestre de actividad regular,
el sistema evita decenas de horas-persona de trabajo manual por ciclo. Al no requerir compra de hardware
especializado ni licencias de software comercial, el proyecto es económicamente viable
con los recursos del equipo académico.

\subsection{Factibilidad legal}

El marco legal del proyecto está determinado por su naturaleza académica, su uso de
herramientas de código abierto y el contexto experimental en que operará: un laboratorio
de farmacología preclínica que trabaja con especímenes de roedores bajo protocolos
aprobados por el IPN.

\subsubsection{Privacidad y manejo de información}

El sistema no registra ni procesa datos personales de seres humanos. Los únicos datos
que almacena son la conducta de cada segundo, duraciones de conducta y metadatos del
experimento (grupo, fecha, tratamiento). Las credenciales de los investigadores
registrados en la plataforma se manejan bajo las directrices de la Ley Federal de
Protección de Datos Personales en Posesión de los Particulares (LFPDPPP), empleando
almacenamiento cifrado y acceso autenticado con \textit{JSON Web Token} (JWT).

\subsubsection{Uso ético de especímenes}

Los procedimientos experimentales con especímenes son realizados íntegramente por el personal
del Laboratorio de Bioquímica Estructural de la ENMyH-IPN. El equipo de desarrollo
recibe únicamente los archivos de video ya grabados y no participa en ningún procedimiento
con los especímenes. El laboratorio opera en apego a la NOM-062-ZOO-1999
\cite{nom062}, norma que define los criterios técnicos, sanitarios y éticos para el
manejo de especímenes de laboratorio en investigación biomédica.

\subsubsection{Licencias de software y propiedad del código}

Cada herramienta del \textit{stack} tecnológico se distribuye bajo una licencia que
permite su uso, modificación y distribución en contextos académicos sin costo:
PSF (Python), BSD-3-Clause (Flask, scikit-learn y PyTorch), MIT (React), PostgreSQL
License (PostgreSQL) y Apache~2.0 (Docker y OpenCV).
El código fuente, los modelos entrenados y la documentación
producidos durante el TT son obra intelectual del equipo desarrollador y quedan
amparados por la Ley Federal del Derecho de Autor (LFDA).

\subsubsection{Marco institucional del IPN}

El proyecto se inscribe en el programa de Trabajo Terminal de la ESCOM-IPN y se rige
por los Lineamientos de Propiedad Intelectual y Transferencia de Conocimiento del
IPN (2023). Su aplicación es exclusivamente científica y educativa, sin
aprovechamiento comercial. Los videos del protocolo FST son propiedad del grupo de
investigación de la ENMyH-IPN; el equipo desarrollador accede a ellos en virtud de la
colaboración académica formal establecida entre los directores del TT y el Dr.\
César Augusto Sandino Reyes López.

% -----------------------------------------------------------
\section{Análisis de sostenibilidad}
% -----------------------------------------------------------

\subsection{Sostenibilidad ambiental}

Reducir el análisis manual quita la necesidad de que los investigadores vayan al
laboratorio solo para cronometrar videos. Al ser web, el sistema se puede usar desde
cualquier lugar. La política de borrar videos a los 30~días evita que el disco crezca
sin control. Según las directrices de la norma ISO~14064-1 \cite{iso14064}, el consumo
energético de procesar un video de 5~min en una computadora local es marginal comparado
con las horas de trabajo que reemplaza.

\subsection{Sostenibilidad social e institucional}

El sistema está hecho para el laboratorio de la ENMyH, pero su diseño modular permite
adaptarlo a otros paradigmas conductuales (laberinto en cruz elevada, campo abierto,
etc.) reconfigurando solo el módulo de clasificación. Eso abre la puerta a que otros
grupos del IPN lo usen sin desarrollarlo desde cero.

Si funciona bien, los investigadores del laboratorio dejaran de cronometrar videos a mano. Eso es lo que
el laboratorio necesita; cualquier beneficio más amplio ---reproducibilidad entre
grupos, escalabilidad a otros paradigmas--- depende de que primero funcione para
el caso para el que fue construido.

\subsection{Mantenibilidad del sistema}

El sistema se entrega como prototipo funcional. Para mantenerlo hace falta conocimiento
básico de Python y Docker. La documentación técnica de este documento y el
\texttt{README} del repositorio están diseñados para que cualquiera que retome el
proyecto pueda entenderlo sin el equipo original. Las herramientas principales (OpenCV,
PyTorch, Flask, React, PostgreSQL) tienen desarrollo activo y comunidades grandes.